\documentclass[11pt]{article}
\usepackage{notes}
\usepackage{subfiles}

\begin{document}

% ---------- course title page ----------
\begin{titlepage}
\centering
\vspace*{3cm}
{\Huge\bfseries\color{coursedark} PHYS-E0510 -- Advanced Computational Methods in Physics D\par}
\vspace{0.8cm}
{\Large\color{courseaccent} Lecture Notes\par}
\vspace{0.3cm}
{\large Sessions 1 to 5\par}
\vspace{2cm}
{\large Jose Lado\par}
\vspace{0.5cm}
{\large Department of Applied Physics, School of Science\\Aalto University, Finland\par}
\vspace{0.5cm}
{\large 2026\par}
\vfill
{\small Each session's title box below names its companion Jupyter notebook and
prerequisites.\par}
\end{titlepage}

\section*{How this course is organized}
\addcontentsline{toc}{section}{\textbf{How this course is organized}}

Every session of this course comes in five pieces, and they are meant to be used together rather
than one instead of another. The recording is the lecture itself, with the derivations done on the
fly and the questions that come up along the way. The slides are the sequence of definitions,
figures and equations in the order they were presented, which makes them the quickest way to find
your place in a recording or to go over a session before the exam. These lecture notes follow the
same sessions in the same order, but write the physics out in connected prose, with the steps a
slide compresses into one line carried through by hand, so you can read them instead of watching a
session a second time, or next to the slides to fill in what a bullet point skips. The Jupyter
notebook of each session, named in the box under the session title and pointed at again in every
\emph{Computational companion} box inside the text, is where the physics is actually computed:
nothing there is a fixed figure, and you are meant to run the cells, change the parameters and see
what happens. Finally, the exercise shorts
(\url{https://joselado.github.io/Advanced_Computational_Methods_Physics_2026/exercise_shorts/}) are
four short silent films per session, one to two minutes each, where we see one idea of the session
with every number on screen computed with the same codes as the notebooks.

None of these pieces is complete on its own: the recordings and slides give the narrative and the
intuition, the notes give derivations you can check line by line, and the notebooks give the
numerics, which build an intuition that none of the others can. A good way to work through a
session is to watch the recording once, read the matching notes with the notebook open next to
them, and then go back to the notebook and change things until the results stop feeling like
something that happened on someone else's screen.

\subsection*{The computational libraries}

All the numerics of the course are done with two Python libraries, rather than with code written
from scratch inside the notebooks, and both have to be installed before the notebooks run; the first
cell of each notebook says which environment variables they expect.

\begin{itemize}[nosep]
  \item \textbf{\texttt{pyqula}} (\url{https://github.com/joselado/pyqula}) does single-particle and
    mean-field electronic structure: tight-binding Hamiltonians, band structures, and the
    self-consistent Hubbard and charge-density-wave calculations of Session 1, whose notebook,
    \texttt{classical\_magnetism.ipynb}, is built on it.
  \item \textbf{\texttt{dmrgpy}} (\url{https://github.com/joselado/dmrgpy}) does the quantum
    many-body calculations: exact diagonalization of spin and fermionic Hamiltonians, and the
    matrix-product-state and DMRG calculations of the later sessions, so the notebooks of Sessions 2
    to 5 are built on it.
\end{itemize}

Both libraries are used here as tools: you do not need to read their source to follow the course,
only to call the functions the notebooks show you.

\clearpage
\tableofcontents
\clearpage

\phantomsection\addcontentsline{toc}{section}{\textbf{Session 1: Classical Magnetism and Mean-Field Theory}}
\subfile{classical_magnetism}

\clearpage

\phantomsection\addcontentsline{toc}{section}{\textbf{Session 2: Exact Diagonalization of the Heisenberg Spin Chain}}
\subfile{quantum_magnetism}

\clearpage

\phantomsection\addcontentsline{toc}{section}{\textbf{Session 3: Exact Diagonalization in the Fermionic Basis}}
\subfile{many_body_fermions}

\clearpage

\phantomsection\addcontentsline{toc}{section}{\textbf{Session 4: Matrix Product States for Spin Chains}}
\subfile{mps_spin_chains}

\clearpage

\phantomsection\addcontentsline{toc}{section}{\textbf{Session 5: Matrix Product States for Fermionic Models}}
\subfile{mps_fermionic}

\end{document}
